% unit U2
\section{What Report~1 Provides}\label{sec:recap}
\mathtgt

This section recalls the pieces of \Rone that Sections~\ref{sec:target} and~\ref{sec:source} use and says what changes from the miniature. Section~\ref{sec:target-canon} recalls the one theorem the proof reuses, with its hypotheses.

\paragraph{Types and evidence.}
An object type of \FCdot is a telescope of propositions about a self, and $\top$ is the empty one, $\obj{\tnil}$ (Section~\ref{R1-sec:mini-types} of report~1). A block name $\sel{x}{\ell}$ names the member $\ell$ of a variable $x$, whatever the type of $x$. The rule for $\ev{member}$ reads one proposition of an object type at an atom (Figure~\ref{R1-fig:evidence} of report~1).

\paragraph{Atoms and roots.}
There is no subsumption in \FCdot. Where DOT retypes a variable, \FCdot wraps it in a cast, and a variable under casts is an atom (Section~\ref{R1-sec:mini-terms} of report~1). These atoms are not the capture atoms of the miniature (Section~\ref{sec:mini-roots}), the elements of a capture set. The full calculus adds atoms that fold and unfold an object type at its variable, and an atom that pairs two typings of one variable (Section~\ref{R1-sec:atoms} of report~1). The root of an atom, $\rt{a}$, is the variable under all of these. Types and capture sets see only the root of an atom. So a cast changes what is known about a variable, never which variable it is. An atom has one root, a term variable. A capture set has a set of roots, which are capture variables, and a variable $x$ has those of $\cset{x}$ (Section~\ref{sec:cc-types}).

\paragraph{Inert stores.}
A context has opaque binders, which a lambda parameter or a let introduces, and transparent binders, which define the block of their variable. A store holds one value per location, a lambda or a literal, never a cast. Store typing $\storej{\sigma}{\Gamma}$ types each value at its precise type, in the context of the entries before it, and binds its location transparently (Section~\ref{R1-sec:mini-machine} of report~1). The store is inert, so the type of a location asserts what its value defines and nothing more.

\paragraph{The machine.}
A state $\mst{\sigma}{K}{t}$ has a store, a continuation of let frames and cast frames, and a term (Figure~\ref{R1-fig:mini-machine} and Figure~\ref{R1-fig:machine} of report~1). Allocation stores $\mathit{core}(v)$, the value under its casts, and continues with $\mathit{adjust}(u, v)$, which hands the stripped casts to the body $u$. Application through a variable is $\beta$ on the stored lambda. Application through a cast atom normalizes the chain of casts on the atom. The head form it finds says how to cast the argument and the result.

\paragraph{Canonical forms.}
Report~1's Theorem~\ref{R1-thm:canon} normalizes closed evidence over a typed store to a typed head form, by structural induction on its typing (Section~\ref{R1-sec:idea} of report~1). Progress inverts the head form by its shape, never a typing derivation of evidence.

\paragraph{Safety by translation.}
Report~1's Theorem~\ref{R1-thm:dot-safety} runs a DOT program and its translation side by side. Its invariant is that some typed target state has the erasure of the source state, and the backward erasure simulation and preservation keep it. Progress of the target, reflected back through erasure, shows that the source never gets stuck.

\paragraph{What changes.}
Section~\ref{sec:mini} split a type into a shape $S$ and a capture set $C$. In the full calculus the shapes are report~1's types and a box shape (Section~\ref{sec:target-boxes}). Evidence for an inclusion of types pairs shape evidence with capture evidence (Section~\ref{sec:target-delta}). Shape evidence keeps an earlier version of report~1's head forms (Section~\ref{sec:mech-trees}). Capture evidence gets none, as Section~\ref{sec:mini} showed. Capture members bring bad bounds back for capture sets (Section~\ref{sec:problem-bad}).
